% Generated draft: requires mathematical and source review.
\documentclass[11pt]{article}
\usepackage[margin=25mm]{geometry}
\usepackage{fontspec}
\setmainfont{Latin Modern Roman}
\usepackage{amsmath,amssymb,mathtools,unicode-math}
\setmathfont{Latin Modern Math}
\usepackage{enumitem,needspace,xurl,hyperref}
\setlength{\parindent}{0pt}
\setlength{\parskip}{5pt}
\setlength{\emergencystretch}{3em}
\newcommand{\N}{\mathbb{N}}
\newcommand{\Z}{\mathbb{Z}}
\newcommand{\Q}{\mathbb{Q}}
\newcommand{\R}{\mathbb{R}}
\newcommand{\C}{\mathbb{C}}
\newcommand{\F}{\mathbb{F}}
\newcommand{\A}{\mathbb{A}}
\newcommand{\PP}{\mathbb{P}}
\newcommand{\E}{\mathbb{E}}
\newcommand{\eps}{\varepsilon}
\newcommand{\dd}{\,\mathrm{d}}
\begin{document}
\section*{UnsolvedMath ID: 800258}
\textbf{Source problem code:} catalogue entry
\subsection{Definitions and mathematical statement}
Let \(n\in\Z\) with \(n\geq 2\). A domain \(D\subset\C^n\) is a nonempty connected open set. It is bounded if it is contained in some Euclidean ball. Under the usual convention, smooth boundary means that \(\partial D\) is a \(C^\infty\) real hypersurface; the supplied record uses the word ``smooth'' without explicitly specifying a differentiability class, so this convention requires confirmation against the source. A holomorphic self-map \(f:D\to D\) is proper if \(f^{-1}(K)\) is compact in \(D\) for every compact subset \(K\subset D\). It is biholomorphic if it is bijective and its inverse \(f^{-1}:D\to D\) is holomorphic.

The question is whether, for every integer \(n\geq 2\), every bounded domain \(D\subset\C^n\) with smooth boundary, and every proper holomorphic map \(f:D\to D\), the map \(f\) is biholomorphic. No additional geometric hypothesis on \(D\) is stated in the supplied record.
\subsection{Short English statement}
Must every proper holomorphic map from a smooth bounded domain in \(\C^n\), with \(n>1\), onto itself be a biholomorphism? The precise boundary regularity intended by ``smooth'' should be confirmed.
\subsection{Sources}
\begin{itemize}
\item Local UnsolvedMath registry, permanent ID \(800258\), source code \texttt{catalogue entry}, title ``Several Complex Variables Problem List --- Problem 57.'' The supplied record attributes the problem to Buzzard, Fornaess, and Krantz and gives the year 1995.
\item Buzzard, Fornaess, and Krantz, \emph{The Several Complex Variables Problem List} (1995), Problem 57. Supplied source URL: \url{https://arxiv.org/abs/math/9509202}. The supplied provenance records access on 2026-10-16 and extraction by \texttt{arxiv\_latex}; these are supplied provenance claims, not independently verified access or extraction.
\item The local record labels the problem ``open.'' This is a supplied status label, not an independently verified assessment of the mathematical literature.
\end{itemize}
\end{document}
